\section{Introduction}
\label{sec:intro}

Prediction markets rank among the most accurate forecasting mechanisms known to the research community. By allowing participants to trade contracts whose payoffs depend on future events, these markets harness collective intelligence to produce probability estimates that routinely outperform individual experts, polls, and statistical models~\cite{arrow2008promise,wolfers2004prediction}. Polymarket, the largest decentralized prediction market, sustains over 600 simultaneous active markets covering domains as diverse as cryptocurrency price movements, geopolitical conflicts, award ceremonies, macroeconomic indicators, and professional sports outcomes. Daily trading volume reaches tens of millions of dollars, and the platform's real-time odds have become reference signals for journalists, analysts, and policymakers alike.

Yet despite the demonstrated power of crowd-aggregated predictions, a persistent gap remains between the information available to individual forecasters and the speed at which markets incorporate that information. Consider a concrete example: when predicting whether a particular film will win Best Picture at the Academy Awards, the market price reflects the aggregate belief of thousands of traders. However, many of these traders may not have systematically analyzed the film's performance across 37 precursor award ceremonies, each of which carries different predictive weight for Oscar outcomes. A forecaster who automatically retrieves, structures, and reasons over these precursor signals can identify windows where the market has not yet fully absorbed available information---and in our experiments, this approach improved prediction accuracy from 71\% to 82\% on a market trading at 74\%.

This example illustrates a broader phenomenon: large language models possess strong reasoning capabilities but lack access to current information, while prediction markets have access to crowd wisdom but may be slow to incorporate domain-specific context. Bridging this gap requires not just a better model, but a better \emph{process}---one that automatically identifies what information is missing, retrieves it from heterogeneous sources, and structures it for effective reasoning.

Existing approaches to LLM-based forecasting fall short in several ways. Zero-shot prompting produces predictions that are poorly calibrated and systematically biased, particularly for events involving current prices or recent developments that postdate the model's training data~\cite{halawi2024approaching}. Retrieval-augmented generation (RAG) improves access to current information but applies a generic retrieval strategy regardless of domain, missing domain-specific signals such as implied volatility for commodity events or precursor correlations for award predictions. Fine-tuning on historical prediction data conflates memorization of past events with genuine forecasting ability, raising concerns about evaluation validity on temporally adjacent events.

We propose Murmur, a system built on three interlocking mechanisms. The \textbf{Context Planner} addresses the information gap by classifying each prediction event into one of eleven domains, identifying knowledge gaps specific to that domain, generating targeted search queries, and structuring retrieved context using Tetlock's superforecaster methodology~\cite{tetlock2015superforecasting}. The \textbf{Evolution Engine} treats the prediction strategy as source code rather than a fixed model, applying LLM-driven code mutations against a fitness function defined by Brier Score on resolved events. Unlike parameter tuning, this approach can discover qualitatively new analysis patterns---for instance, learning to apply contrarian adjustments specifically to mid-range market prices where crowd overconfidence is empirically most prevalent. The \textbf{Virtual Prediction Market} extends evaluation beyond historical backtesting by simulating an exchange populated with heterogeneous trading agents (noise traders, informed traders, momentum followers, contrarians, and market makers), enabling strategy robustness testing under dynamic market conditions that static backtesting cannot capture.

Our work makes the following contributions:

\begin{itemize}
    \item We introduce the Context Planner, a domain-aware pipeline that automatically detects knowledge gaps and injects structured real-time context before LLM-based prediction. Across 11 domains, context enrichment yields 8--11\% accuracy improvement in information-sensitive categories while providing negligible benefit in efficiently-priced markets, confirming the domain-dependent nature of prediction alpha.
    
    \item We propose evolutionary code search for prediction strategy discovery, demonstrating that LLM-driven code mutation with single-change constraints and diversity pressure can match market-level Brier Scores. We document critical design decisions---warm-start initialization, banned modification tracking, and constrained prompting---that distinguish successful evolution from naive approaches that consistently fail near local optima.
    
    \item We present an empirical analysis of prediction alpha across 11 Polymarket domains, identifying the conditions under which LLM-augmented forecasting provides genuine edge: low-liquidity markets, events with delayed information incorporation, and domains where structured decomposition outperforms intuitive crowd judgment.
    
    \item We demonstrate architectural generality by showing that Murmur's core mechanisms---context-aware agent reasoning, evolutionary strategy discovery, and agent-based simulation---transfer to supply chain optimization and healthcare coordination, suggesting a unified framework for decision-making under uncertainty.
\end{itemize}
